\section{Conclusion}

This survey studies recent progress on using LLMs to strengthen OS security by proposing a component-based taxonomy, \revise{namely}, kernel, drivers, networking, file systems, system services, and hypervisors/TEEs.
Across these components, \revise{recurring approaches use LLMs as} \emph{semantic integrators} that translate heterogeneous artifacts (code, specifications, logs/traces, crash reports, and tool outputs) into actionable hypotheses, and as \emph{tool-oriented agents} that improve the efficiency of workflows such as specification formalization, stateful fuzzing, and patching in generate--validate loops.
At the same time, OS security imposes stringent correctness and trust requirements, so LLM outputs must be treated as candidates and grounded in deterministic evidence (e.g., reproducers, symbolized traces, diffs, and tests), with reviewable artifacts and guardrails that preserve correctness, compatibility, and performance. 
By highlighting both \revise{recurring} design patterns and their failure modes, our survey clarifies \revise{where evaluated applications exist and where evidence remains limited} (e.g., hypervisor/TEE and some driver-facing tasks), and which future directions are most urgent---especially evidence-grounded evaluation and review-centric pipelines that scale human governance rather than overwhelm it.




\begin{comment}
This survey reviews recent progress on applying LLMs to OS security and organizes the literature around major OS components. In the introduction and background, we motivate this area through the lens of trust boundaries and argue that many practical attacks and defenses are shaped by how untrusted inputs reach privileged code and data. Building on this organization, we summarize how recent work uses LLMs to assist common security workflows, including evidence triage, vulnerability reasoning, test and patch assistance, fuzzing support, and security policy understanding. We also discuss the risks that arise when language models are embedded in privileged pipelines, and we highlight safeguards based on least privilege tool access, isolation, and comprehensive auditing. A key requirement for dependable deployment is to couple model suggestions with independent validation when possible, using deterministic checks and machine checkable evidence such as tests, code diffs, traces, or verification tools. Looking forward, progress will benefit from clearer evaluation that focuses on security outcomes, as well as benchmarks with maintainable ground truth and validation centered workflows with safe fallback behavior. We hope this survey helps guide the development of OS security systems that use language models effectively while remaining auditable and safe to deploy.
\zmy{LLMs augment OS security workflows → but introduce privileged-path risks → requiring validation, auditing, and evaluation.}
\end{comment}
